% Appendix
\appendix

\chapter{Project Repository Structure}
\label{app:project}

The project repository is organized as follows:

\section*{Directory Structure}
\begin{verbatim}
Project/
├── analysis/           # Jupyter notebooks for analysis
├── backup/             # Backup data files
├── dataset/            # Raw datasets
├── models/             # Model implementation modules
├── Report/             # LaTeX report files
├── results/            # Output files and processed data
├── scripts/            # Utility scripts
├── requirements.txt    # Python dependencies
└── README.md           # Project documentation
\end{verbatim}

\section*{Analysis Notebooks}
\begin{itemize}
    \item \texttt{fire.ipynb}: Fire dataset exploration
    \item \texttt{climate.ipynb}: Climate data analysis
    \item \texttt{elevation.ipynb}: DEM processing
    \item \texttt{landcover.ipynb}: Land cover analysis
    \item \texttt{soil.ipynb}: Soil data processing
    \item \texttt{merge.ipynb}: Dataset integration
    \item \texttt{feature\_engineering.ipynb}: Feature creation
    \item \texttt{sampling.ipynb}: Negative sample generation
    \item \texttt{knn.ipynb}, \texttt{decision\_tree.ipynb}, \texttt{random\_forest.ipynb}: Supervised models
    \item \texttt{kmeans.ipynb}, \texttt{dbscan.ipynb}: Clustering analysis
\end{itemize}

\section*{Data Sources}
\begin{itemize}
    \item \textbf{Fire data}: NASA FIRMS (MODIS/VIIRS)
    \item \textbf{Land cover}: FAO Global Land Cover Share
    \item \textbf{Climate}: CHIRPS precipitation, TerraClimate temperature
    \item \textbf{Elevation}: GMTED2010 DEM
    \item \textbf{Soil}: HWSD v2.0
\end{itemize}

\chapter{Preprocessing Details}
\label{app:preprocessing}

\section*{Fire Dataset Statistics}
\begin{table}[H]
\centering
\caption{Fire Dataset Quality Summary}
\begin{tabular}{|l|c|c|}
\hline
\textbf{Attribute} & \textbf{Missing (\%)} & \textbf{Action} \\
\hline
Latitude/Longitude & 0.1 & Removed \\
Brightness & 0.0 & None \\
FRP & 2.3 & Median imputation \\
Confidence & 1.8 & Mode imputation \\
Acquisition date & 0.0 & None \\
\hline
\end{tabular}
\end{table}

\section*{Derived Topographic Variables}

\textbf{Slope calculation} using Horn's algorithm:
\begin{equation}
    \text{Slope} = \arctan\left(\sqrt{\left(\frac{\partial z}{\partial x}\right)^2 + \left(\frac{\partial z}{\partial y}\right)^2}\right)
\end{equation}

\textbf{Aspect calculation}:
\begin{equation}
    \text{Aspect} = \arctan2\left(-\frac{\partial z}{\partial y}, -\frac{\partial z}{\partial x}\right)
\end{equation}

\textbf{Terrain complexity}: Standard deviation of elevation in a 3×3 pixel window.

\textbf{Coastal proximity}: Euclidean distance to nearest coastline boundary in kilometers.

\section*{Climate Feature Aggregation}

Seasonal aggregations computed:
\begin{itemize}
    \item \texttt{prec\_winter}: December-February mean precipitation
    \item \texttt{prec\_summer}: June-August mean precipitation
    \item \texttt{tmax\_summer}: June-August mean maximum temperature
    \item \texttt{tmin\_winter}: December-February mean minimum temperature
    \item \texttt{total\_annual\_precip}: Sum of monthly precipitation
    \item \texttt{annual\_temp\_range}: Annual maximum minus minimum temperature
\end{itemize}

\textbf{Drought Index} formula:
\begin{equation}
    \text{DI} = \frac{T_{summer} - \bar{T}_{summer}}{\sigma_T} - \frac{P_{summer} - \bar{P}_{summer}}{\sigma_P}
\end{equation}

Higher values indicate drier, hotter conditions.

\chapter{Hyperparameter Tuning Details}
\label{app:hyperparams}

\section*{KNN Hyperparameters}
\begin{table}[H]
\centering
\caption{KNN Cross-Validation Results}
\begin{tabular}{|c|c|c|}
\hline
\textbf{k} & \textbf{Mean F1} & \textbf{Std} \\
\hline
3 & -- & -- \\
5 & -- & -- \\
7 & -- & -- \\
9 & -- & -- \\
11 & -- & -- \\
15 & -- & -- \\
21 & -- & -- \\
\hline
\end{tabular}
\end{table}

\section*{Decision Tree Hyperparameters}
Grid search space:
\begin{itemize}
    \item \texttt{max\_depth}: \{3, 5, 7, 10, 15, 20, None\}
    \item \texttt{min\_samples\_split}: \{2, 5, 10, 20\}
    \item \texttt{min\_samples\_leaf}: \{1, 2, 5, 10\}
\end{itemize}

\section*{Random Forest Hyperparameters}
Grid search space:
\begin{itemize}
    \item \texttt{n\_estimators}: \{50, 100, 200, 300\}
    \item \texttt{max\_depth}: \{10, 15, 20, None\}
    \item \texttt{min\_samples\_split}: \{2, 5, 10\}
    \item \texttt{max\_features}: \{sqrt, log2, 0.3\}
\end{itemize}

Best parameters found: [To be filled from experiments]

\section*{Complete Feature List}
\begin{table}[H]
\centering
\caption{All Features Used in Modeling}
\begin{tabular}{|l|l|}
\hline
\textbf{Category} & \textbf{Features} \\
\hline
Geographic & latitude, longitude, lat\_long\_interaction \\
\hline
Climate - Temperature & tmax\_summer, tmax\_winter, tmin\_summer, tmin\_winter, \\
 & annual\_temp\_range, summer\_heat\_index \\
\hline
Climate - Precipitation & prec\_summer, prec\_winter, prec\_autumn, prec\_spring, \\
 & total\_annual\_precip, summer\_precip\_ratio \\
\hline
Climate - Derived & drought\_index \\
\hline
Topography & elevation, slope, aspect\_sin, aspect\_cos, \\
 & terrain\_complexity, coastal\_proximity, \\
 & elevation\_slope\_interaction \\
\hline
Soil & ORG\_CARBON, TOTAL\_N, PH\_WATER, SAND, \\
 & CLAY, SILT, fine\_fraction, ph\_carbon\_interaction \\
\hline
\end{tabular}
\end{table}

\chapter{Clustering Details}
\label{app:clustering}

\section*{K-Means Configuration}
\begin{itemize}
    \item Initialization: k-means++ (intelligent seeding)
    \item Maximum iterations: 300
    \item Number of initializations: 10 (best result selected)
    \item Convergence tolerance: $1 \times 10^{-4}$
\end{itemize}

\section*{K-Means Evaluation Metrics}
\begin{table}[H]
\centering
\caption{K-Means Clustering Quality}
\begin{tabular}{|l|c|}
\hline
\textbf{Metric} & \textbf{Value} \\
\hline
Number of Clusters (k) & -- \\
Inertia (WCSS) & -- \\
Silhouette Score & -- \\
Calinski-Harabasz Index & -- \\
Davies-Bouldin Index & -- \\
\hline
\end{tabular}
\end{table}

\section*{DBSCAN Configuration}
\begin{itemize}
    \item Epsilon ($\epsilon$): Selected from k-distance knee point
    \item Min\_samples: dimensionality + 1
    \item Metric: Euclidean distance
\end{itemize}

\section*{DBSCAN Results}
\begin{table}[H]
\centering
\caption{DBSCAN Clustering Results}
\begin{tabular}{|l|c|}
\hline
\textbf{Metric} & \textbf{Value} \\
\hline
Number of Clusters & -- \\
Number of Noise Points & -- \\
Percentage Noise & -- \\
Silhouette Score (excl. noise) & -- \\
\hline
\end{tabular}
\end{table}

\section*{CLARANS Configuration}
\begin{itemize}
    \item Number of clusters: Same as K-Means
    \item numlocal: 2
    \item maxneighbor: max(250, 1.25\% × k × (n-k))
\end{itemize}

\section*{Clustering Comparison Summary}
\begin{table}[H]
\centering
\caption{Clustering Performance Comparison}
\begin{tabular}{|l|c|c|c|}
\hline
\textbf{Metric} & \textbf{K-Means} & \textbf{DBSCAN} & \textbf{CLARANS} \\
\hline
Clusters Found & -- & -- & -- \\
Silhouette Score & -- & -- & -- \\
Computation Time (s) & -- & -- & -- \\
\hline
\end{tabular}
\end{table}

\chapter{Additional Results}
\label{app:results}

\section*{Confusion Matrices}

\textbf{Random Forest Confusion Matrix:}
\begin{table}[H]
\centering
\begin{tabular}{|c|c|c|}
\hline
 & \textbf{Predicted No Fire} & \textbf{Predicted Fire} \\
\hline
\textbf{Actual No Fire} & 1766 (TN) & 19 (FP) \\
\hline
\textbf{Actual Fire} & 85 (FN) & 680 (TP) \\
\hline
\end{tabular}
\end{table}

\section*{Feature Importance (Complete)}
\begin{table}[H]
\centering
\caption{All Feature Importance Scores}
\begin{tabular}{|l|c|}
\hline
\textbf{Feature} & \textbf{Importance} \\
\hline
coastal\_proximity & 0.114 \\
lat\_long\_interaction & 0.067 \\
elevation & 0.060 \\
drought\_index & 0.058 \\
prec\_winter & 0.054 \\
elevation\_slope\_interaction & 0.044 \\
tmax\_summer & 0.043 \\
prec\_autumn & 0.040 \\
total\_annual\_precip & 0.039 \\
terrain\_complexity & 0.039 \\
\hline
\multicolumn{2}{|c|}{[Additional features...]} \\
\hline
\end{tabular}
\end{table}

\section*{Python Dependencies}
\begin{verbatim}
numpy>=1.21.0
pandas>=1.3.0
scikit-learn>=1.0.0
geopandas>=0.10.0
rasterio>=1.2.0
matplotlib>=3.4.0
seaborn>=0.11.0
jupyter>=1.0.0
\end{verbatim}
